% Ad Familiares 13.15 (ad-familiares-13-15)
% Source: https://raw.githubusercontent.com/PerseusDL/canonical-latinLit/master/data/phi0474/phi056/phi0474.phi056.perseus-lat1.xml
% Fetched by scripts/fetch_latin.py

% Dateline (Perseus): Scr. Romae ex. m. Mad. a. 709 (45) .

\ciceroLetterOpener{CICERO CAESARI IMP. S.}

\ciceroSection{1} Precilium tibi commendo unice, tui necessari, mei familiarissimi, viri optimi filium. quem cum adulescentem ipsum propter eius modestiam, humanitatem, animum et amorem erga me singularem mirifice diligo, tum patrem eius re doctus intellexi et didici mihi fuisse semper amicissimum. em hic ille est de illis , maxime qui inridere atque obiurgare me solitus est, quod me non tecum, praesertim cum abs te honorificentissime invitarer, coniungerem ; a)ll' e)mo\n ou)/ pote qumo\n e)ni\ sth/qessin e)/peiqen. audiebam enim nostros proceres clamitantis : a)/lkimos e)/ss', i(/na ti/s se kai\ o)yigo/nwn e)u\ ei)/ph|. w(\s fa/to, to\n d' a)/xeos nefe/lh e)ka/luye me/laina.

\ciceroSection{2} sed tamen idem me consolantur etiam ; hominem perustum etiamnum gloria volunt incendere atque ita loquuntur : mh\ ma\n a)spoudi/ ge kai\ a)kleiw=s a)poloi/mhn, a)lla\ me/ga r(e/cas ti kai\ e)ssome/noisi puqe/sqai. sed me minus iam movent, ut vides. itaque ab Homeri magniloquentia confero me ad vera praecepta *eu)ripi/dou : misw= sofisth/n, o(/stis ou)x au(tw=| sofo/s, quem versum senex Precilius laudat egregie et ait posse eundem et a(/ma pro/ssw kai\ o)pi/ssw videre et tamen nihilo minus ai)e\n a)risteu/ein kai\ u(pei/roxon e)/mmenai a)/llwn.

\ciceroSection{3} sed ut redeam ad id unde coepi, vehementer mihi gratum feceris, si hunc adulescentem humanitate tua, quae est singularis, comprehenderis et ad id, quod ipsorum Preciliorum causa te velle arbitror, addideris cumulum commendationis meae. genere novo sum litterarum ad te usus, ut intellegeres non vulgarem esse commendationem.
